\chapter*{Introduction}
\label{cap:introduction}
\addcontentsline{toc}{chapter}{Introduction}

The computing landscape is undergoing a major shift. Driven by the rapid growth of
Artificial Intelligence, real-time data processing is no longer merely desirable; it
has become a fundamental requirement.
While the traditional approach relied on offloading neural network computation to
large cloud data centers, current trends are moving toward Edge AI: bringing
intelligence closer to the data source, directly onto the local device.
This approach significantly reduces latency, improves user privacy by avoiding the
transmission of sensitive data over the network, and reduces the load on communication
infrastructures. \citep{zhou2019edgeintelligencepavingmile}.

This is where FPGAs become a key enabling technology.
Unlike traditional CPUs and GPUs, which process data either sequentially or through
fixed forms of parallelism, FPGAs allow custom hardware architectures to be designed
from the ground up, optimizing data flow and exploiting instruction-level parallelism.
This is particularly relevant for addressing current physical limitations such as the
Power Wall, where further increasing the operating frequency of conventional processors
is no longer energy-efficient.
\citep{10.1145/3282307}

For these reasons, this Bachelor's Thesis focuses on an in-depth exploration of
hardware-software co-design.
To carry this out, the Digilent Zybo platform will be used \citep{digilentZyboOverview},
a development board that integrates the flexibility of an ARM Cortex-A9 processor for
control-oriented tasks with the computational capabilities of a Xilinx/AMD 7-series
FPGA for intensive processing workloads, all within the same chip. From an academic
perspective, this platform provides an ideal environment for emulating industrial
embedded systems, giving the designer fine-grained control over every clock cycle and
every data transfer through high-performance buses such as AXI4.

\section*{Motivation}
\addcontentsline{toc}{section}{Motivation}
The demand for computational performance, mainly driven by Artificial Intelligence and
Deep Learning algorithms, is increasing at a pace that general-purpose processors can
no longer address efficiently under strict power-consumption constraints. Moving
inference from the cloud to local embedded devices is no longer optional; it is the
natural direction of the field.

The main technical challenge of this work, and what makes it particularly demanding,
is to make the most of an FPGA with limited resources, only around 80 DSP blocks, in
order to run Machine Learning models. This will be supported by modern High-Level
Synthesis (HLS) tools, which help reduce the complexity traditionally associated with
pure hardware design and make the development process more accessible and efficient.

\section*{Objectives}
\addcontentsline{toc}{section}{Objectives}
The main objective of this work is to design, implement, and evaluate a heterogeneous
embedded system based on the Zynq-7010 SoC, capable of accelerating the inference of a
lightweight Artificial Intelligence model.
As a direct consequence, this work also aims to quantify the performance improvement
achieved by the hardware implementation when compared with the execution of the same
inference flow on the ARM microprocessor.
To achieve this goal, the following specific objectives are defined:

\begin{enumerate}
  \item \textbf{Platform bring-up:} Configure the Vivado/Vitis development environment
    from scratch for the Digilent Zybo board, validating the independent operation of
    both the ARM Cortex-A9 microprocessor and the FPGA.
  \item \textbf{Communication infrastructure design:} Implement high-bandwidth
    communication channels between the Processing System and the Programmable Logic,
    using the AMBA AXI protocol and Direct Memory Access (DMA) controllers.
  \item \textbf{Acceleration of a baseline computation:} Develop a basic linear algebra
    IP core in C++, specifically a matrix multiplication accelerator, and synthesize it
    into hardware using HLS in order to validate the data flow between the processor
    and the FPGA.
  \item \textbf{AI implementation under limited resources:} Adapt a neural network model,
    specifically a CIFAR-10 Convolutional Neural Network, to the limited resources of
    the Zynq-7010, generating its hardware implementation and developing the control
    software required to provide the input data.
  \item \textbf{Performance evaluation (Benchmarking):} Compare execution times, in terms
    of latency, by running the neural network inference on the CPU alone and comparing
    it with the FPGA-assisted execution.
\end{enumerate}

\section*{Work plan}
\addcontentsline{toc}{section}{Work plan}
To organize the development process and meet the project milestones, the work is divided
into the following incremental phases:

\textbf{Phase 1.} Research and familiarization: Review of the state of the art in Zynq-based
systems, AXI buses, and High-Level Synthesis. Installation of the required tools
(Vivado and Vitis) and execution of a "Hello World" application on the board.

\textbf{Phase 2.} PS-PL communication: Configuration of control registers through AXI-Lite and
setup of the memory subsystem using AXI DMA. Development of bare-metal C drivers to
transfer data from RAM to the FPGA.

\textbf{Phase 3.} Development of the baseline accelerator in HLS: Implementation of matrix
multiplication in C/C++. Use of HLS pragmas such as PIPELINE and UNROLL to optimize
the generated hardware. Integration of the resulting IP block into the Vivado design
together with the DMA.

\textbf{Phase 4.} Design and implementation of the AI model: Training and/or export of a
Convolutional Neural Network model for CIFAR-10 image classification. Generation of
the neural network hardware accelerator and integration into the Vivado block design.

\textbf{Phase 5.} Testing and benchmarking: Execution of inference on the FPGA and comparison
against the hls4ml-generated core running on the ARM. Collection of performance metrics
and verification of the obtained results.

\textbf{Phase 6.} Thesis report writing: Documentation of the development process, analysis of
the obtained results, and final writing of the Bachelor's Thesis document.
